\section{Formal Foundations}
    \label{sec:foundations}
    
    \subsection{Primitive Structure of the Theory}
        
        We define the Swirl-String Theory (SST) as a continuum-based framework built on a minimal set of primitive quantities.
        
        \textbf{Primitive set:}
        \begin{align}
            \mathcal{P} = \{ \rho_f, v_{\swirlarrow}, \omega_c \}
        \end{align}
        
        where:
        \begin{itemize}
            \item $\rho_f$ is the effective background fluid density
            \item $v_{\swirlarrow}$ is the characteristic swirl velocity
            \item $\omega_c$ is the Compton angular frequency
        \end{itemize}
        
        \textbf{[ORTHODOX]}: $\omega_c = \frac{m_e c^2}{\hbar}$ is a standard physical constant.
        
        \textbf{[SPECULATIVE]}: The identification of $\omega_c$ as a primitive dynamical frequency of the vortex core is an SST postulate.
    
    \subsection{Derived Length Scale}
        
        The characteristic core radius is defined as:
        
        \begin{align}
            r_c = \frac{v_{\swirlarrow}}{\omega_c}
            \label{eq:rc_definition}
        \end{align}
        
        \textbf{[DERIVED]}: This removes $r_c$ as an independent parameter and eliminates circular definitions.
    
    \subsection{Circulation Quantization}
        
        The circulation is defined as:
        
        \begin{align}
            \Gamma = \oint \mathbf{v} \cdot d\ell
        \end{align}
        
        We introduce the fundamental circulation quantum:
        
        \begin{align}
            \Gamma_0 = 2\pi r_c v_{\swirlarrow}
            \label{eq:gamma0}
        \end{align}
        
        and impose:
        
        \begin{align}
            \Gamma = n \Gamma_0, \quad n \in \mathbb{Z}
        \end{align}
        
        \textbf{[ORTHODOX]}: Circulation quantization appears in superfluid systems.
        
        \textbf{[DERIVED]}: In SST, $\Gamma_0$ is no longer an input but follows from Eq.~\ref{eq:rc_definition}.
    
    \subsection{Energy and Mass Relation}
        
        Mass is interpreted as stored rotational energy:
        
        \begin{align}
            E = m c^2
        \end{align}
        
        \textbf{[ORTHODOX]}: Standard relativistic energy relation.
        
        \textbf{[SPECULATIVE]}: SST interprets this energy as localized vortex kinetic energy.
    
    \subsection{Swirl Velocity Constraint}
        
        We define the dimensionless ratio:
        
        \begin{align}
            \eta_0 = \frac{v_{\swirlarrow}}{c}
        \end{align}
        
        Empirically:
        
        \begin{align}
            \eta_0 \approx \frac{\alpha}{2}
        \end{align}
        
        \textbf{[DERIVED]}: This connects the swirl velocity to the fine-structure constant.
    
    \subsection{Swirl-Clock Definition}
        
        We introduce the Swirl-Clock as a local scalar field encoding relativistic time dilation:
        
        \begin{align}
            S_{(t)} = \sqrt{1 - \frac{v^2}{c^2}}
            \label{eq:swirl_clock}
        \end{align}
        
        \textbf{[ORTHODOX]}: This matches the Lorentz factor.
        
        \textbf{[DERIVED]}: In SST, it is interpreted as a local fluid-kinematic effect.
    
    \subsection{Density Hierarchy}
        
        We distinguish three densities:
        
        \begin{align}
            \rho_f &\quad \text{(background fluid)} \\
            \rho_E &= \frac{E}{c^2} \\
            \rho_{core} &\quad \text{(saturated vortex core)}
        \end{align}
    
    \subsection{Core Density Closure}
        
        The core density follows from a force constraint:
        
        \begin{align}
            \rho_{core} \sim \frac{m_e c^2}{2\pi v_{\swirlarrow}^2 r_c^3}
            \label{eq:core_density}
        \end{align}
        
        \textbf{[DERIVED]}: This eliminates $\rho_{core}$ as an independent parameter.
    
    \subsection{Geometric Representation}
        
        A swirl-string is represented as a closed curve:
        
        \begin{align}
            X(t,\sigma) = (x(t,\sigma), y(t,\sigma), z(t,\sigma))
        \end{align}
        
        with $\sigma \in [0,2\pi)$.
    
    \subsection{Summary of Dependency Structure}
        
        The full dependency chain is:
        
        \begin{align}
        (\rho_f, v_{\swirlarrow}, \omega_c)
            \rightarrow r_c
            \rightarrow \Gamma_0
            \rightarrow \rho_{core}
        \end{align}
        
        \textbf{[CRITICAL]}:
        No quantity is defined circularly.
    
    \subsection{Consistency Statement}
        
        \begin{itemize}
            \item All derived quantities depend only on $\mathcal{P}$
            \item No external parameters are introduced
            \item The theory is closed under its primitive set
        \end{itemize}
    
    \subsection{Interpretation}
        
        The Formal Foundations establish SST as:
        
        \begin{itemize}
            \item a continuum theory with minimal primitives
            \item a geometrically grounded framework
            \item a non-circular parameter system
        \end{itemize}